% Flowchart for the Methodology section.
%
% Drawn with TikZ so no binary asset is needed and the fonts match the paper. A two-column IEEE
% layout gives about 3.4in (~86mm) of column width, so node text is set at \scriptsize and every
% node is kept well inside that: an earlier version used 63mm text widths plus side-positioned
% labels and overflowed the column by 17pt.
%
% Self-contained: this file carries its own figure environment, caption and label, so main.tex
% needs only a single \input.
\begin{figure}[t]
\centering
\definecolor{flowA}{RGB}{235,242,250}
\definecolor{flowB}{RGB}{244,238,228}
\definecolor{flowEdge}{RGB}{90,100,115}

\begin{tikzpicture}[
  font=\scriptsize,
  >={Stealth[length=1.5mm]},
  node distance=2.4mm,
  box/.style={draw=flowEdge, rounded corners=1pt, align=center, inner sep=2.2pt,
              minimum height=5mm},
  prob/.style={box, fill=flowA, text width=52mm},
  meth/.style={box, fill=flowB, text width=17mm},
  wide/.style={box, fill=flowB, text width=58mm},
  ar/.style={->, draw=flowEdge, line width=0.4pt},
]

% 1. the problem, with its closed-form answer key
\node[prob] (prob) {\textbf{Linear--quadratic HJB}, $p$ assets:\\
                    value function known in \textbf{closed form}};

% 2. four solver families, side by side
\node[meth, below=4mm of prob, xshift=-28.5mm] (a) {finite\\differ-\\ences};
\node[meth, right=1.6mm of a]                   (b) {Howard\\policy\\iteration};
\node[meth, right=1.6mm of b]                   (c) {direct\\PINN\\~};
\node[meth, right=1.6mm of c]                   (d) {policy-\\iteration\\PINN};

% 3. the two ways of obtaining a control
\node[wide, below=6mm of b, xshift=-9mm] (grid)
  {\textbf{grid:} $V$ at $N^p$ nodes $\rightarrow$ \textbf{reconstruct} $\hat V$
   $\rightarrow$ differentiate};
\node[wide, below=2.2mm of grid] (net)
  {\textbf{network:} $V$ is continuous $\rightarrow$ \textbf{autograd}
   differentiates it exactly};

\draw[ar] (prob.south) -- ++(0,-1.6mm) -| (a.north);
\draw[ar] (a.south) |- ([yshift=1.2mm]grid.north west);
\draw[ar] (b.south) |- ([yshift=1.2mm]grid.north);
\draw[ar] (c.south) |- ([yshift=1.2mm]net.north);
\draw[ar] (d.south) |- ([yshift=1.2mm]net.north east);

% 4. the protocol, applied to every method
\node[wide, below=6mm of net] (proto)
  {\textbf{score \emph{value} and \emph{control}} against the closed form,\\
   on identical states, at matched accuracy};
\draw[ar] (grid.south) -- ++(0,-1.6mm) -| ([xshift=-2mm]proto.north);
\draw[ar] (net.south)  -- ++(0,-1.6mm) -| ([xshift=2mm]proto.north);

% 5. what the protocol reveals
\node[wide, below=4mm of proto] (find)
  {value accuracy does not bound control accuracy; the control error is
   governed by the \textbf{reconstruction order}};
\draw[ar] (proto.south) -- (find.north);

\end{tikzpicture}
\caption{The benchmark design. A problem whose answer is known in closed form is solved by four
methods that differ in \emph{how they obtain a control}: the two grid methods must reconstruct a
continuous function from nodal values and then differentiate it, while the two network methods are
continuous functions by construction and are differentiated exactly. Every solver is scored on both
the value function and the control, against the analytic solution and never against another
solver.}
\label{fig:method}
\end{figure}
